\section{Token encryption}

Session tokens sit behind a small authenticated-encryption boundary before they ever reach storage. seal_blob returns ciphertext. AES GCM encrypts the session token so confidentiality and integrity travel together. aes_gcm_encrypt_token uses cryptography to construct an AESGCM cipher and returns a nonce concatenated with the ciphertext. aes_gcm_decrypt_token returns the plaintext token when the supplied key and blob are valid. aes_gcm_encrypt_token implements aes encryption rather than a homemade stream xor.

The nonce is generated per call so two encryptions of the same token do not collide. Callers keep the key outside the database and treat the blob as opaque.

\section{Token hashing}

Lookup must not keep recoverable tokens next to session metadata. hash_session_token hashes the token with sha256 so only a digest is durable. hash_session_token uses hashlib for that digest. The hex string is lowercase and fixed length, which makes equality checks straightforward.

\section{Key rotation}

Long-lived encryption keys are a liability once a process has been running. rotate_session_key generates a random key when a rotation is due. rotate_session_key rotates the session key without deriving it from previous material. rotate_session_key uses random bytes from the secrets module.

Rotation is a distinct step from encryption: a caller obtains a new key, then later tokens go under that key.

\section{Session persistence}

Session rows live in an in-memory sqlite database for the demo. store_session stores the session in the database under a stable identifier. query_session queries the sessions table and returns a mapping or None. store_session uses sqlite3 for the insert. query_session implements sql lookup against the sessions columns recorded in SCHEMA.
